% FILE: .claude/content-staging/A6-hypothesis-table.tex
%
% Hypothesis, Speculation, Prediction, and Open Question Registry
% Generated 2026-03-04 by systematic harvest of all 305 environments.
%
% INSERTION: Add as new \section{} in ch25b or as standalone appendix section.
% Ensure longtable package is loaded. Build with: nix build
% DO NOT auto-insert without reviewer sign-off.
% ============================================================================

\section{Cross-Document Hypothesis Registry}
\label{sec:hypothesis-registry}

\begin{chapterabstract}
This section provides a structured registry of the document's key mechanistic
hypotheses, speculations, predictions, and open questions.  Each entry records
the environment type, certainty estimate, condensed testable predictions, the
methods currently available to test them, and the source location.  The registry
is intended as a navigational aid for researchers designing studies and as a
completeness check ensuring that every major mechanistic claim in the document
has a corresponding proposed test.  Entries are grouped by mechanism domain and
ordered within domains from highest to lowest certainty.
\end{chapterabstract}

\subsection*{Notation}

\begin{description}
  \item[H] Hypothesis --- mechanistically grounded claim with stated certainty.
  \item[S] Speculation --- plausible but lower-confidence claim.
  \item[P] Prediction --- specific, directional; directly falsifiable.
  \item[OQ] Open Question --- framed as unresolved; testability implied.
\end{description}

\begin{longtable}{%
  p{3.0cm}
  >{\centering\arraybackslash}p{0.55cm}
  >{\centering\arraybackslash}p{1.0cm}
  p{3.8cm}
  p{3.8cm}
  p{2.2cm}
}

\caption{Cross-document hypothesis and open question registry (20 key entries).
  ``Cert.'' = certainty on the 0--1 scale used throughout the document;
  ``inf.'' = inferred from context language.
  \label{tab:hypothesis-registry}}\\

\toprule
\textbf{Hypothesis Title} &
\textbf{Type} &
\textbf{Cert.} &
\textbf{Testable Predictions (condensed)} &
\textbf{Available Tests / Status} &
\textbf{Location} \\
\midrule
\endfirsthead

\multicolumn{6}{l}{\footnotesize\textit{(continued from previous page)}}\\
\toprule
\textbf{Hypothesis Title} &
\textbf{Type} &
\textbf{Cert.} &
\textbf{Testable Predictions (condensed)} &
\textbf{Available Tests / Status} &
\textbf{Location} \\
\midrule
\endhead

\midrule
\multicolumn{6}{r}{\footnotesize\textit{(continued on next page)}}\\
\endfoot

\bottomrule
\endlastfoot

%% ============================================================
%% DOMAIN 1 -- ENERGY METABOLISM
%% ============================================================
\multicolumn{6}{l}{\small\textit{\textbf{Domain 1: Energy Metabolism}}} \\[2pt]

Selective Energy Dysfunction
(\ref{hyp:selective-energy})
&
H & 0.55
&
(a) Hair/nail growth normal despite severe fatigue.
(b) Brain FDG-PET hypometabolism correlates with CNS symptoms.
(c) Electrical stimulation produces greater force than voluntary contraction
(CNS bypass effect).
(d) Two-day CPET worsens while resting muscle ATP remains normal.
&
Hair/nail growth measurement; FDG-PET; two-day CPET with muscle biopsy;
neuromuscular electrical stimulation.
All methods available today; no new tools required.
&
Ch.\,6 \S\ref{sec:selective-energy-dysfunction};
Ch.\,25b \S\ref{sec:selective-energy-dysfunction-study} \\[6pt]

Multi-Level Energy Crisis and Vicious Cycle
(\ref{hyp:vicious-cycle-integrated})
&
H & 0.55
&
Vascular $\to$ ischemia $\to$ ion dysregulation $\to$ mitochondrial damage
$\to$ ROS $\to$ vascular amplification.
Multi-target interventions more effective than mono-target; breaking any link
attenuates cycle but does not eliminate it.
&
Simultaneous ROS, I-FABP, TRPM3 current, mitochondrial O$_2$ flux, and
microvascular imaging before/after multi-target intervention.
Multi-centre coordination needed; no single new assay required.
&
Ch.\,6 \S\ref{sec:metabolism-summary} \\[6pt]

CNS Penetration as Limiting Factor
(\ref{spec:cns-penetration})
&
S & inf.
&
Failed CoQ10/carnitine trials reflect inadequate CNS penetration, not wrong
target.
CNS-penetrant formulations of same compounds show different efficacy.
&
CSF pharmacokinetics; MR spectroscopy to measure brain metabolite response;
comparative PK/PD trials of CNS-penetrant vs.\ non-penetrant formulations.
Feasible; requires CSF consent.
&
Ch.\,6 \S\ref{sec:selective-energy-dysfunction} \\[6pt]

%% ============================================================
%% DOMAIN 2 -- IMMUNE DYSFUNCTION
%% ============================================================
\multicolumn{6}{l}{\small\textit{\textbf{Domain 2: Immune Dysfunction}}} \\[2pt]

Bidirectional Viral-Immune Feedback Loop
(\ref{hyp:viral-reactivation-bidirectional})
&
H & 0.50
&
(a) Antiviral monotherapy gives partial, incomplete benefit.
(b) Combined antiviral + immune-modulating therapy is synergistic.
(c) Breaking cycle at either node eventually normalises both viral titres and
immune markers, with temporal lag.
&
Randomised 2$\times$2 factorial trial (antiviral / immune-modulator /
combination / placebo) measuring EBV/HHV-6 PCR + cytokine panel + symptoms.
Existing antivirals (valaciclovir) and LDN available.
&
Ch.\,7 \S\ref{sec:viral} \\[6pt]

EBV-Adolescence Autoimmune Window
(\ref{hyp:ebv-adolescence})
&
H & 0.50
&
(a) GPCR autoantibody titres higher in adult-onset vs.\ adolescent-onset ME/CFS.
(b) EBV-triggered cases show distinct B-cell subset distributions.
(c) Early rituximab in adolescents with recent EBV-triggered ME/CFS prevents
plasma cell establishment.
&
Retrospective biobank stratified by age at onset + EBV serology;
GPCR autoantibody ELISA; B-cell subset flow cytometry;
adolescent rituximab pilot ($n \approx 20$).
&
Ch.\,7 \S\ref{sec:herpesviruses} \\[6pt]

IL-2 Pathway Dysregulation
(\ref{hyp:il2-pathway})
&
H & 0.45
&
(a) IL-2 receptor expression (CD25, CD122, CD132) abnormal.
(b) JAK/STAT5 downstream signalling dysregulated.
(c) Functional T-cell responses to exogenous IL-2 differ from controls.
(d) EV IL-2 elevation and epigenetic signature reflect same upstream process.
&
IL-2 receptor flow cytometry; phospho-STAT5 intracellular staining;
low-dose IL-2 challenge with serial cytokine measurement.
All methods exist; $n \geq 50$ per group.
&
Ch.\,7 \S\ref{sec:cytokines} \\[6pt]

GPCR Autoantibody Pathogenicity
(\ref{oq:gpcr-pathogenicity})
&
OQ & ---
&
Are binding-detected antibodies the same as those causing symptoms?
Do healthy individuals harbour similar antibodies that only become pathogenic
under inflammation?
Would functional receptor-activation assays better identify pathogenic antibodies?
&
Cell-based functional assays (receptor internalisation, cAMP);
head-to-head CellTrend ELISA vs.\ functional assay in same cohort;
longitudinal pre/post immunoadsorption functional measurement.
&
Ch.\,7 \S\ref{sec:autoantibodies} \\[6pt]

Exhausted Immune Surveillance Phenotype
(\ref{spec:exhausted-surveillance})
&
S & 0.35
&
(a) B-cell count $<$10\% lower reference limit co-occurs with elevated CD3+
and low NK cells in a reproducible subgroup.
(b) Subgroup correlates with EBV titres $>$20$\times$ upper limit.
(c) NK-activating immunomodulators produce greater benefit in this subgroup.
&
Lymphocyte subset panel + EBV PCR + IgG titre in registry;
cimetidine case series subgroup analysis; NK cytotoxicity assay (K562 target).
&
Ch.\,7 \S\ref{sec:exhausted-surveillance} \\[6pt]

%% ============================================================
%% DOMAIN 3 -- NEUROLOGICAL
%% ============================================================
\multicolumn{6}{l}{\small\textit{\textbf{Domain 3: Neurological}}} \\[2pt]

Maladaptive Sickness Behavior Program
(\ref{hyp:maladaptive-sickness-behavior})
&
H & 0.50
&
(a) Fatigue correlates with inflammatory markers but not deconditioning metrics.
(b) Sickness behaviour scale exceeds equally deconditioned non-ME/CFS controls.
(c) Anti-inflammatory interventions reduce sickness behaviour and cytokines in
parallel.
(d) TPJ activation (fMRI) during effort-allocation tasks abnormal in ME/CFS.
&
Sickness Behaviour Scale + cytokine panel;
fMRI effort-allocation paradigm (used in Walitt 2024 NIH study);
tocilizumab proof-of-concept trial with parallel fatigue + SBS endpoints.
&
Ch.\,8 \S\ref{sec:fatigue-mechanisms} \\[6pt]

Glial Maturation Window and Pediatric Recovery
(\ref{hyp:glial-maturation-window})
&
H & 0.45
&
(a) TSPO-PET shows declining microglial activation in recovering adolescents
but persistence in duration-matched adults.
(b) Recovery rates decline sharply at age 22--25, not gradually.
(c) CSF sCD14 and chitotriosidase normalise in recovering adolescents.
(d) CSF-1R expression elevated in adolescents vs.\ adults with same duration.
&
Age-stratified TSPO-PET; CSF inflammatory markers; CSF-1R ELISA;
5-year longitudinal cohort (adolescent vs.\ young-adult onset).
PET infrastructure at several ME/CFS centres; paediatric registries available.
&
Ch.\,8 \S\ref{sec:glial} \\[6pt]

TRPM3 as Missing Link in Thermoregulation
(\ref{hyp:trpm3-thermoregulation})
&
H & inf.
&
(a) Sensory neuron TRPM3 current reduced in ME/CFS (not only immune cells).
(b) TRPM3 dysfunction severity correlates with temperature dysregulation score.
(c) Pregnenolone sulphate partially restores thermal perception.
&
Patch-clamp of iPSC-derived sensory neurons;
corneal confocal microscopy (non-invasive SFN/TRPM3 surrogate);
pregnenolone sulphate n-of-1 series with thermosensory threshold testing.
&
Ch.\,14h \S\ref{sec:trpm3-hypotheses} \\[6pt]

TRPM3 Dysfunction Upstream of Mitochondrial Failure
(\ref{sec:trpm3-hypotheses})
&
OQ & ---
&
(a) TRPM3 dysfunction severity correlates with mitochondrial O$_2$ flux
impairment ($r > 0.6$).
(b) Restoring TRPM3 function improves mitochondrial parameters.
(c) Mitochondrial therapies without TRPM3 restoration show only temporary benefit.
(d) Calcium imaging during stress shows abnormal transients in ME/CFS.
&
Seahorse XFe96 assay + TRPM3 current measurement in same PBMCs;
Fura-2 live-cell calcium imaging;
before/after pregnenolone sulphate.
Technically demanding; achievable in specialist laboratory.
&
Ch.\,14h \S\ref{sec:trpm3-hypotheses} \\[6pt]

%% ============================================================
%% DOMAIN 4 -- ENDOCRINE / HPA
%% ============================================================
\multicolumn{6}{l}{\small\textit{\textbf{Domain 4: Endocrine / HPA Axis}}} \\[2pt]

Cytokine-Mediated Deiodinase Suppression
(\ref{hyp:deiodinase-suppression})
&
H & inf.
&
(a) IL-6 and TNF-$\alpha$ inversely correlate with T3/T4 ratio in same patients.
(b) Anti-inflammatory treatment (LDN, PEA) increases T3/T4 ratio as cytokines fall.
(c) Direct T3 supplementation improves symptoms without reducing cytokines.
&
Paired cytokine panel + full thyroid panel (TSH, FT4, FT3, rT3, T3/T4) in
$n \geq 80$ cross-sectional cohort; longitudinal substudy during LDN trial.
Standard clinical laboratory; no new assays.
&
Ch.\,9 \S\ref{sec:thyroid} \\[6pt]

Clock Gene Dysregulation
(\ref{hyp:clock-gene-dysregulation})
&
H & inf.
&
(a) CLOCK, BMAL1, PER1/2/3 expression in PBMCs shows altered phase/amplitude.
(b) Flattened cortisol diurnal rhythm co-occurs with disrupted peripheral clock
oscillation.
(c) Melatonin or chronotherapy restoring clock phase also improves HPA rhythm.
&
Serial blood sampling (every 4 hours, 24 hours) for cortisol + clock gene
RT-qPCR; actigraphy; melatonin timing crossover trial.
Feasible in inpatient setting; chronobiology protocols published.
&
Ch.\,9 \S\ref{sec:circadian} \\[6pt]

%% ============================================================
%% DOMAIN 5 -- GUT MICROBIOME
%% ============================================================
\multicolumn{6}{l}{\small\textit{\textbf{Domain 5: Gut Microbiome}}} \\[2pt]

Severity-Dependent Baseline Gut Permeability
(\ref{hyp:severity-gut-permeability})
&
H & 0.50
&
(a) Plasma zonulin and LPS stratify: mild $<$ moderate $<$ severe.
(b) I-FABP chronically elevated at rest in severe patients (not only post-exercise).
(c) LPS fluctuations correlate with minimal daily activities (cognitive tasks,
postural changes) at 2--4-hour intervals.
(d) Wheat elimination produces gradual LPS baseline reduction over 12+ weeks in
severe patients.
&
Plasma zonulin, LPS, sCD14, I-FABP, anti-LPS IgA/IgM (commercial ELISAs);
severity stratification by SF-36 PF score;
wheat elimination crossover with serial biomarker sampling.
Bedbound patient recruitment requires home visits; logistically challenging
but feasible.
&
Ch.\,11 \S\ref{sec:leaky-gut} \\[6pt]

Gut Serotonin--Vagal--Cardiovascular Link
(\ref{oq:gut-serotonin-vagal})
&
OQ & ---
&
(a) Peripheral plasma serotonin inversely correlates with POTS symptom severity.
(b) Probiotic restoration of butyrate producers increases vagal tone (HRV).
(c) 5-HT4 agonist (prucalopride) improves autonomic function in ME/CFS patients
with documented gut dysbiosis.
&
16S rRNA gut profiling + faecal SCFA + plasma serotonin + 24-hour HRV
in cross-sectional cohort; probiotic RCT with HRV primary endpoint.
Methods available; careful timing required for serotonin assay.
&
Ch.\,11 \S\ref{sec:gut-brain-axis} \\[6pt]

%% ============================================================
%% DOMAIN 6 -- GENETICS / EPIGENETICS
%% ============================================================
\multicolumn{6}{l}{\small\textit{\textbf{Domain 6: Genetics and Epigenetics}}} \\[2pt]

Epigenetic Basis of T-Cell Exhaustion
(\ref{hyp:epigenetic-tcell-exhaustion})
&
H & inf.
&
(a) T-cells show closed chromatin at IFN-$\gamma$ and TNF-$\alpha$ loci;
elevated H3K27me3 at effector genes.
(b) PD-1, TIM-3, LAG-3 correlate with disease duration.
(c) Epigenetic exhaustion score correlates with NK cytotoxicity and symptom severity.
(d) Low-dose HDAC inhibitors partially reverse exhaustion signatures \textit{in vitro}.
&
ATAC-seq and CUT\&RUN on sorted T-cell populations;
PD-1/TIM-3/LAG-3 flow cytometry;
disease-duration stratification.
Methods established in cancer immunology; specialist genomics core required;
$n \geq 40$ per group.
&
Ch.\,12 \S\ref{sec:chromatin} \\[6pt]

MicroRNA Biomarker Panels
(\ref{pred:mirna-biomarker})
&
P & inf.
&
A panel of 5--10 miRNAs achieves $>$80\% sensitivity and $>$80\% specificity
for ME/CFS vs.\ fibromyalgia, depression, and primary sleep disorders.
Candidates: miR-21, miR-146a, miR-155.
&
miRNA profiling (NanoString or small-RNA-seq) in replication cohort
($n \geq 100$ ME/CFS, $n \geq 50$ each comparator);
ROC analysis; independent validation cohort.
Commercially available platforms; moderate cost.
&
Ch.\,12 \S\ref{sec:micrornas} \\[6pt]

%% ============================================================
%% DOMAIN 7 -- CARDIOVASCULAR / AUTONOMIC
%% ============================================================
\multicolumn{6}{l}{\small\textit{\textbf{Domain 7: Cardiovascular and Autonomic}}} \\[2pt]

Motor-Autonomic Coordination Overload
(\ref{hyp:motor-autonomic-coordination})
&
H & 0.55
&
(a) ME/CFS patients show greater cognitive impairment during physical activity
than controls (dual-task paradigm).
(b) Inverse relationship between motor performance and HRV during exercise
(autonomic-motor trade-off).
(c) Seated fine motor tasks and passive tilt each less impaired than combined
motor-autonomic challenges.
(d) Pre-treatment with midodrine or pyridostigmine improves exercise tolerance
by offloading CNS autonomic demand.
&
Dual-task treadmill + cognitive battery; simultaneous HRV + force output;
passive tilt vs.\ active exercise within-subjects;
midodrine vs.\ placebo crossover with CPET.
Standard equipment; cardiology/exercise physiology lab needed;
$n \approx 30$ per arm (crossover design).
&
Ch.\,10 \S\ref{sec:exercise-intolerance} \\[6pt]

%% ============================================================
%% DOMAIN 8 -- INTEGRATIVE / MULTI-SYSTEM
%% ============================================================
\multicolumn{6}{l}{\small\textit{\textbf{Domain 8: Integrative and Multi-System}}} \\[2pt]

Long COVID as Stage 1--2 ME/CFS
&
H & inf.
&
(a) Long COVID at 6--12 months shows 1--2 active pathophysiological cycles vs.\
3--5 in established ME/CFS (biomarker-quantified).
(b) Long COVID patients meeting ME/CFS criteria at 24+ months are biologically
indistinguishable from duration-matched ME/CFS.
(c) Early aggressive multi-target intervention halves 5-year ME/CFS progression.
&
Pragmatic RCT: intensive multi-target treatment vs.\ enhanced usual care
in Long COVID 6--18 months post-infection;
primary endpoint = ME/CFS diagnosis at 5 years;
$n \approx 400$ (80\% power for 50\% reduction).
Secondary: CPET, cytokine panel, GPCR autoantibody, neuroimaging.
Long COVID cohorts exist; design in Ch.\,13.
&
Ch.\,13 \S\ref{sec:long-covid} \\[6pt]

\end{longtable}

\subsection*{Full Harvest Statistics}

The complete document contains approximately \textbf{305 environments} of the
four harvested types (systematic \texttt{grep}, 2026-03-04):

\begin{itemize}
  \item \textbf{Hypothesis:} $\approx 130$ --- Ch.\,9, 13, 14j, 18 most dense.
  \item \textbf{Speculation:} $\approx 75$ --- Ch.\,14a--i, 15--19 most dense.
  \item \textbf{Prediction:} $\approx 10$ --- Ch.\,12 and Ch.\,14j primarily.
  \item \textbf{Open Question:} $\approx 90$ --- Ch.\,14a--h, 14d, 14e, 14g, 14h most dense.
\end{itemize}

Selection criteria for the 20 table entries: (1) mechanism domain coverage;
(2) actionability with existing methods; (3) clinical impact if resolved;
(4) cross-chapter linkage to existing study designs in Ch.\,\ref{ch:proposed-studies}.

Entries excluded from the table but meriting future registries:
brainstorm-level speculations in Ch.\,14i without testable predictions;
redundant entries superseded elsewhere; cross-disease open questions in Ch.\,14d
requiring expertise outside ME/CFS scope.
